%% 毕业设计小结

\begin{designsummary}
本毕业设计以无人机在非定常大气环境中的节能飞行为研究背景，
围绕"基于强化学习的自推进体在非定常流场中的能耗优化迁移研究"这一课题，
完成了从物理建模、算法设计到数值仿真与结果分析的完整研究流程。

在物理建模方面，构建了二维非定常双涡流流场模型，
推导了自推进体的运动学方程，定义了推进功率与总能量消耗的量化方法，
并设计了"时间对齐"的直线基准策略以确保能耗对比的公平性。

在算法设计方面，基于最大熵框架实现了Soft Actor-Critic连续控制算法，
构建了8维状态空间和2维连续动作空间，
设计了以最小能耗和成功到达为核心的复合奖励函数，
完成了全套训练流程的编程实现与调试。

在仿真分析方面，系统研究了智能体迁移策略的轨迹形态、能耗对比、
参数影响、收敛性、泛化能力和物理机制。
结果表明，经SAC算法训练的智能体能够在双涡流流场中自主习得
"借流而行"的迁移策略，其能量消耗仅为直线基准策略的约三分之一，
且策略具备良好的泛化能力。

在理论分析方面，通过余弦相似度指标揭示了智能体在
"顺流省力"与"横向寻路"之间保持约60度最优夹角的物理机制，
使研究超越了纯数值对比，深入到了流体物理与智能控制的交叉领域。

通过本次毕业设计，本人深入掌握了强化学习算法的原理与实现方法，
提升了对活性物质物理和非平衡输运理论的理解，
锻炼了将理论模型转化为数值仿真程序并进行分析验证的能力，
为今后的科研工作奠定了坚实的基础。
\end{designsummary}
